\section{Conclusion}
\label{sec:conclusion}

We studied whether validation-tuned L2 regularization improves held-out generalization for small tabular classification datasets relative to an otherwise matched logistic-regression baseline. Across four datasets and three random seeds per dataset, L2 produced a clear and statistically supported reduction in the train-validation gap, but only a small and non-significant increase in aggregate test accuracy and weighted F1.

The key takeaway is simple: in this setting, L2 regularization is a reliable anti-overfitting mechanism, not a universal accuracy booster. That is a useful distinction for practitioners and for empirical papers that interpret regularization effects too broadly.

Several extensions would strengthen the picture. Future work should repeat the protocol for small multilayer perceptrons, replace single validation splits with repeated nested cross-validation, add calibration and confusion-matrix analysis, and test whether the same pattern persists on a larger collection of small tabular datasets or stratified subsamples of larger benchmarks.
